\scanpage{89}
\mspage{88}

\uline{Remarque} 3.13. Voici une généralisation parfois utile de 3.5. b). Considérons une partie $E$ de $X$ telle que pour tout morphisme étale $f:X'\quad X$, $E'=f^{-1}(E)$ soit ``très dense'' dans $X'$ (EGA IV 10.1.3.)\unsure{;} i.e. que pour tout $f$ comme ci-dessus, tout ouvert $U$ de $X'$ qui contient $E'$ soit égal à $X'$. Alors les foncteurs fibres $F_{\bar x}$ sur $X_{\text{ét}}$, pour $x\quad E$, forment déjà une ``famille conservative'' (et par suite les corollaires 3.6., 3.7. et 3.8. restent valables en se bornant à y choisir $x\quad E$). Cela se voit immédiatement en reprenant la démonstration donnée de 3.5. b). Le résultat précédent s'applique notamment dans les siutations suivantes:

a) $X$ est un préschéma de \unsure{Jawobson|Jacobson} (EGA IV 10.4.1.), par exemple localement de type fini sur un corps, ou sur $\operatorname{Spec}(\underline{Z})$, et $E$ est l'ensemble des points fermés de $X$.

b) $X$ est localement de type fini sur un préschéma $S$, et $E$ est l'ensemble des points de $X$ qui sont fermé dans leur fibre.

c) $X$ est localement noethérien, et $E$ est l'ensemble des $x\quad X$ tels que $\bar x$ soit un ensemble fini (i.e. un schéma semi-local de dimension $\quad 1$) cf EGA IV 10.5.3. et 10.5.5.
